% Source: sources/returns/S6_04_MAPS_MANIFOLDS/sections/manifolds.tex:6339–6585
\paragraph{Theorem D.2.31 (Proved: all zero-degree original caps and their free defects).}
In the natural candidate catalogue we may retain every original
additive cap at which $P$ is narrow, without requiring a single
component orbit. At $I_n^*$ this theorem uses the original cap;
finite good-reduction caps at the other additive sites have all
the choices and twists of D.2.30. The zero-free multiplicative
caps and the unique rank-two zero retain their preceding marked
models. These choices give $228$ additive-filling patterns on
$160$ natural initial rows. Exactly $30$ patterns have a positive
original component defect. The complete data are
\path{records/manifolds_original_zero_degree.json}.

For an original additive cap let $\rho_i$ be the number of
orbits of $Q$ on its components, and put
\[
 D=\sum_{i\text{ original additive}}(\rho_i-1).
\]
The only nonzero local defects in this catalogue are
\[
\begin{array}{c|rrrr}
\text{type}&III^*&IV^*&I_3^*&I_1^*\\\hline
\rho_i&5&3&3&2\\
\rho_i-1&4&2&2&1.
\end{array}                                                   \tag{D.OD1}
\]
Every other permitted original cap has $\rho_i=1$.

Compute the Seifert group and its Smith cokernel
$\mathbb Z^r\oplus A_1$ as before, treating an unchanged cap
as an order-one meridian with zero numerator in the canonical
marking. Let $T$ be the torsion of the matrix in D.2.30,
using its finite and regular columns and setting both fixed
sweeps of every original additive cap to zero. Then the full
integral groups are
\[
\begin{array}{c|c}
j&H_j(X;\mathbb Z)\\\hline
0,6&\mathbb Z\\
1&\mathbb Z^r\oplus A_1\\
2&\mathbb Z^{2r+D}\oplus T\\
3&\mathbb Z^{2r+2D}\oplus T\\
4&\mathbb Z^{2r+D}\oplus A_1\\
5&\mathbb Z^r.
\end{array}                                                   \tag{D.OD2}
\]
The complete fundamental group is D.2.23, including every
order and affine twist. Thus the component defect adds free
homology, with no uncomputed torsion extension.

There is a full smooth identification in the $b=1$ sector.
For every simply connected member, let $M$ be the product of
its regular logarithmic multiplicities. Then
\[
 X\cong
 \mathop{\#}^{D}(S^2\times S^4)
 \ \#\ \mathop{\#}^{D}(S^3\times S^3)
 \ \#\ \widetilde\Sigma_0 N_M,                             \tag{D.OD3}
\]
with the oriented twisted suspension and the spin five-manifold
$N_M$ defined in D.2.14. The allowed free defects in this sector
are $D=0,1,2,4$. The identification retains the full refined
spin data, including even $M$; it is not inferred from
additive homology alone.

\begin{proof}
\emph{1. The actual original cap.}
At an additive fibre the condition that $P$ is narrow means
that $\mathcal O(P-O)$ has degree zero on every component.
Its underlying line on the rational Kodaira tree is trivial.
The unit equivariant multiplier can be removed by its
rigidified logarithm isotopy, exactly as in D.2.26; the tree
has $H^1=0$, and the isotopy has zero meridian winding on the
normal collar. The cap consequently retracts, with its marked
angular circle, onto
\[
 S^1_\delta\times\operatorname{MT}_Q(K),
\]
where $K$ is the underlying Kodaira tree and $Q$ acts by its
actual component translation. The nontrivial contraction
modulus is retained as the mapping-torus direction; only its
unit circle part is straightened. This is the same relative
comparison proved for the odd-star caps in D.2.16, and uses
only narrowness of $P$, not the number of component orbits.

The tree has $H_1(K)=0$ and one free degree-two class per
component. The permutation invariants and coinvariants are
both free of rank $\rho$. The mapping-torus sequence and
Kunneth therefore give
\[
 H_1(U)=\mathbb Z^2,\quad H_2(U)=\mathbb Z^{\rho+1},\quad
 H_3(U)=\mathbb Z^{2\rho},\quad H_4(U)=\mathbb Z^\rho.
                                                               \tag{D.OD4}
\]
There is no permutation torsion in these groups.

The two full normal tori are still primitive. The $\delta$
normal torus bounds the section disc times its circle. For
the other, take the least power $Q^q$ that is narrow and its
local Neron logarithm. Its marked loop extends across the
base disc; the rigidified norm of the unit linearization has
no meridian winding, so it introduces no fibre correction.
The local existence and extension of this logarithm are as
in \cite[Section 2.4, equation (2.2)]{manUU21} and D.2.16.
For each actual matrix the fixed lattice has basis
$\delta,\ell$, with $\gamma^*(\ell)=q$. Thus these two
normal tori project to the \emph{entire} integral fixed
lattice, not a finite-index sublattice. Their values under
the exterior fibre functional are zero in the stated
rigidified marking.

We also retain the component cokernel. Since (D.OD4) is free,
Poincare--Lefschetz duality gives
$H_2(U,\partial U)=H^4(U)=\mathbb Z^\rho$.
The map from $H_2(U)$ is the folded component-intersection
map. On the component classes its rank is $\rho-1$:
the Kodaira intersection matrix is negative semidefinite
with one-dimensional radical generated by the fibre
multiplicity vector, and averaging to permutation invariants
preserves that radical. The exact folded matrices are also
retained in the certificate. The principal fibre relation
bounds the full rank by $\rho-1$, so the additional
$\delta$--mapping-torus class cannot increase it.

It follows from the pair sequence that
\[
 \operatorname{im}\bigl(H_2(\partial U)\to H_2(U)\bigr)
 \text{ is primitive of rank }2,
 \qquad
 H_2(U)/\operatorname{im}H_2(\partial U)
                 \cong\mathbb Z^{\rho-1}.                 \tag{D.OD5}
\]
Primitivity is integral: the quotient is the image of a map
into the free group $H_2(U,\partial U)$.
The boundary $H_2$ is free of rank four in every such marked
additive case. Its kernel has rank two, and the two primitive
fixed-loop normal tori just constructed therefore exhaust it.
Indeed their projection is already the whole fixed lattice,
and the fibre part is injective rationally. This proves that
the only change in the local Mayer--Vietoris data is the
free summand in (D.OD5).

At a zero-free multiplicative site in this natural catalogue,
$P$ narrow implies that $Q$ is transitive on its components.
The direct component data and the primitive vanishing
matrices verify this at every occurrence. Hence an allowed
zero-degree original multiplicative cap has $\rho=1$ and
is already covered by the earlier primitive comparison.
The additive component permutations give exactly (D.OD1).

\emph{2. Global integral homology.}
The subgroups in (D.OD5) are primitive, so choose integral
splittings in all the original cap groups. Eliminating their
boundary images in Mayer--Vietoris leaves their direct sum
$\mathbb Z^D$, together with exactly the torsion presentation
of D.2.30. The other cap boundary maps are onto as proved
there. The original caps have the same first Betti number as
their finite or minimal comparisons, so the degree-one bridge
kernel is still free of rank $r$. Consequently
$H_2=\mathbb Z^{2r+D}\oplus T$.
Every added original cap has Euler number zero, and
$\chi(X)=2$ remains unchanged. Duality gives all of (D.OD2),
including $b_3=2r+2D$ and the paired middle torsion.
This proves the formula without an implicit integral
spectral-sequence splitting.

There are $396$ exact local zero-degree original certificates.
Enumerating ``original'' precisely when $P$ is narrow,
and ``finite log'' at every finite good-reduction site,
with an original cap required at $I_n^*$ in the present
sector, gives the stated $228$ patterns on $160$ rows.
No logarithmic order is bounded in this enumeration.

\emph{3. Smooth classification when $b=1$.}
For a simply connected member all remaining finite logs have
primitive actual gamma numerator. Since $b=1$, their indices
$q_i$ are one; a site of fixed-gamma index greater than one
would contradict simple connectivity. D.2.30 and the
regular-log calculation therefore give
\[
 H_2(X)=\mathbb Z^D\oplus(\mathbb Z/M)^2,
 \qquad H_3(X)=\mathbb Z^{2D}\oplus(\mathbb Z/M)^2.
\]
The scalar circle is free on every original cap, by the
product description in part 1, and free on every finite
or regular logarithmic cap by the primitive gamma numerator.
At each zero-free multiplicative cap, $Q=G$ is transitive
and the pair consisting of its actual primitive vanishing
vector and $\delta$ has unit Smith factors. The program
checks this pair at every such site in the relevant patterns.
The rank-two zero has the same full original short lattice
and the same unique fixed sphere as D.2.14, with normal
bundle $\mathcal O(-1)_{+1}\oplus\mathcal O(-1)_{-1}$.
Thus this is the actual semifree circle action on $X$.

The fixed sphere is nullhomotopic. Its specialization comes
from the common fibre class $s$; it has no contribution to
the extra component cokernels in (D.OD5). The submatrix on
$s$ and the delta sweeps is the unimodular Seifert matrix
when $b=1$ and $\pi_1(X)=1$. It therefore kills $s$
integrally. Hurewicz gives the assertion about the sphere.

We check spin despite the new free second homology. On every
original cap the holomorphic volume form is the product of
a nonvanishing canonical two-form on the local minimal
elliptic surface with $dv/v$. Translation by $Q$ preserves
the first factor, and the change of $dv/v$ is a base-surface
one-form whose wedge with that two-form is zero. The volume
form descends, so the canonical line of that cap is trivial,
as in D.2.16. The free part of $H_2(X)$ is generated by the
component cokernels of these caps. Hence $c_1(X)$ evaluates
to zero on it. Since $X$ is simply connected,
$H^2(X;\mathbb Z)$ is free, so $c_1(X)=0$ and $X$ is spin.

The equivariant surgery of D.2.14 now gives
$Y=X\#(S^3\times S^3)$ with a free circle action and a
projectable spin structure. Its orbit five-manifold is
simply connected and spin, and the homotopy/Gysin sequences
give
\[
 (B,e)\cong
 \left(\mathop{\#}^{D+1}(S^2\times S^3)\#N_M,
                        (x,0,\ldots,0)\right).
\]
Here $e$ is primitive; the spin Smale--Barden classification
and its realization theorem identify this pair, as in
\cite[Theorems 2.2--2.3]{manBarden65}.
Apply the spin branch of \cite[Theorem A]{manGR25}
successively, retaining the zero Euler class on the other
summands. Together with its suspension formula this gives
\[
 Y\cong\mathop{\#}^{D}(S^2\times S^4)
       \#\mathop{\#}^{D+1}(S^3\times S^3)
       \#\widetilde\Sigma_0N_M.
\]
The full Zhubr classification permits the same cancellation
of one $S^3\times S^3$ summand as in D.2.14 and D.2.16;
it retains the refined class $p$ and the full class
$\mu\in H_6(K(H_2,2);\mathbb Z)$, including all free cubic,
mixed and finite data. The added $S^3\times S^3$ maps
constantly to this Eilenberg--Mac Lane space and spin-bounds
$S^3\times D^4$. It therefore changes neither bordism
invariant and adds only two to $b_3$. The diffeomorphism of
the stabilized manifolds identifies these complete data,
which classify the unstabilized manifolds as well.
This proves (D.OD3), including its oriented framing and
all two-primary data.
\end{proof}
